\documentclass[10pt,aspectratio=169]{beamer}
\usepackage[default]{lato}
\usepackage{amsmath,amssymb,booktabs,array}
\definecolor{blue}{RGB}{0,114,178}
\definecolor{red}{RGB}{213,94,0}
\definecolor{yellow}{RGB}{240,228,66}
\setbeamercolor{structure}{fg=blue}
\setbeamercolor{frametitle}{fg=blue}
\setbeamercolor{section in head/foot}{fg=blue}
\setbeamertemplate{navigation symbols}{}
\setbeamersize{text margin left=0.7cm,text margin right=0.7cm}
\useinnertheme[shadow]{rounded}
\setbeamertemplate{itemize items}{-}
\setbeamercolor{block title}{fg=blue,bg=blue!20!white}
\setbeamercolor{block body}{bg=blue!5!white}
\setbeamertemplate{footline}{\hbox{\begin{beamercolorbox}[wd=\paperwidth,ht=3ex,dp=1.5ex,leftskip=2ex,rightskip=2ex]{page footer}\insertsectionnavigationhorizontal{0.7\paperwidth}{}{} \hfill \insertframenumber{} / \inserttotalframenumber\end{beamercolorbox}}}
\newcommand{\E}{\mathbb{E}}
\newcommand{\Var}{\operatorname{Var}}
\newcommand{\Cov}{\operatorname{Cov}}
\newcommand{\Prb}{\operatorname{Pr}}
\title[]{\textcolor{blue}{Study Guide: Lectures 1--3\\Probability, Statistical Inference, and OLS}}
\author{25117 - Econometrics}
\institute[shortinst]{Universitat Pompeu Fabra}
\date{}
\begin{document}

\begin{frame}{Study Guide: Lectures 1--3}
\begin{center}
{\Large\textcolor{blue}{Probability $\longrightarrow$ Inference $\longrightarrow$ Regression}}\\[.4cm]
25117 -- Econometrics \quad Universitat Pompeu Fabra
\end{center}
\begin{itemize}
\item \textbf{Lecture 1:} describe distributions and dependence; understand random sampling, estimators, and their sampling distributions.
\item \textbf{Lecture 2:} quantify uncertainty with standard errors, tests, confidence intervals, and comparisons of means.
\item \textbf{Lecture 3:} derive and interpret ordinary least squares, assess fit, and distinguish statistical association from causality.
\end{itemize}
\begin{block}{How to study}
For each formula, identify the object being estimated, the meaning of each symbol, the assumptions, and the interpretation. Work the examples before checking the answers.
\end{block}
\end{frame}

\section{Lecture 1}
\begin{frame}{Parameters, Random Variables, Estimators, and Estimates}
\begin{center}
\begin{tabular}{lll}
\toprule
Object & Notation & What is random?\\
\midrule
Population mean & $\mu_Y=\E[Y]$ & Fixed, usually unknown\\
Random sample & $Y_1,\ldots,Y_n$ & Values vary across samples\\
Observed sample & $y_1,\ldots,y_n$ & Recorded numbers\\
Estimator of the mean & $\overline Y=n^{-1}\sum_iY_i$ & A function of random data\\
Estimate of the mean & $\overline y=n^{-1}\sum_i y_i$ & One realized number\\
\bottomrule
\end{tabular}
\end{center}
\begin{itemize}
\item Capitals denote random variables; lower-case letters denote their realizations.
\item Example: observing $(y_1,y_2,y_3)=(170,180,175)$ gives $\overline y=175$ cm.
\item The \textbf{sampling distribution} describes the different values of $\overline Y$ across repeated samples of the same size.
\end{itemize}
Unknown does not mean random: in frequentist inference, $\mu_Y$ is fixed even when we do not know it.
\end{frame}

\begin{frame}{Probability Distributions: Mass, Density, and Cumulative Probability}
\small
A random variable assigns a number to each outcome of an experiment. An \textbf{event} is a set of outcomes.
\begin{columns}[T]
\begin{column}{.48\textwidth}
\begin{block}{Discrete: probability mass function}
\[
p_Y(y)=\Prb(Y=y),\qquad \sum_y p_Y(y)=1.
\]
Example: for a Bernoulli variable,
\[
\Prb(Y=1)=p,\qquad \Prb(Y=0)=1-p.
\]
Each possible value can have positive probability.
\end{block}
\end{column}
\begin{column}{.48\textwidth}
\begin{block}{Continuous: probability density function}
\[
f_Y(y)\geq0,\qquad \int_{-\infty}^{\infty}f_Y(y)\,dy=1.
\]
\[
\Prb(a<Y\leq b)=\int_a^b f_Y(y)\,dy.
\]
Probability is \textbf{area}, not height; $\Prb(Y=y)=0$ at any single point.
\end{block}
\end{column}
\end{columns}
\medskip
For either type, the \textbf{cumulative distribution function (CDF)} is
\[
F_Y(y)=\Prb(Y\leq y),\qquad
\Prb(a<Y\leq b)=F_Y(b)-F_Y(a).
\]
\end{frame}

\begin{frame}{Essential Distinction: Standard Deviation and Standard Error}
\small
For independent and identically distributed observations with variance $\sigma_Y^2$:
\begin{center}
\begin{tabular}{p{3.3cm}p{4.3cm}p{5.6cm}}
\toprule
Quantity & Formula & Interpretation\\
\midrule
Population standard deviation & $\sigma_Y=\sqrt{\Var(Y_i)}$ & Spread of individual observations\\[5pt]
SD of the sample mean & $\sigma_{\overline Y}=\sigma_Y/\sqrt n$ & Spread of means across repeated samples\\[5pt]
Sample standard deviation & $s_Y=\sqrt{\frac{\sum_i(y_i-\overline y)^2}{n-1}}$ & Estimate of $\sigma_Y$ from one sample\\[9pt]
Estimated standard error & $\displaystyle SE(\overline Y)=\frac{s_Y}{\sqrt n}$ & Estimated sampling uncertainty of the mean\\
\bottomrule
\end{tabular}
\end{center}
\begin{block}{When do we replace $\sigma_Y$ with $s_Y$?}
When $\sigma_Y$ is unknown, as it usually is. We still divide by $\sqrt n$:
$s_Y$ alone is \emph{not} the standard error of $\overline Y$.
\end{block}
Example: if $n=100$ and $s_Y=20$ cm, then $SE(\overline Y)=2$ cm.
More observations make the mean more precise; they need not make individuals more similar.
\end{frame}

\begin{frame}{Moments: Mean, Variance, and Standard Deviation}
\begin{align*}
\mu_Y=\E[Y]&=\sum_y y\,p_Y(y)
\quad\text{or}\quad \int_{-\infty}^{\infty}y f_Y(y)\,dy,\\
\sigma_Y^2=\Var(Y)&=\E[(Y-\mu_Y)^2]=\E[Y^2]-\mu_Y^2,\\
\sigma_Y&=\sqrt{\sigma_Y^2}.
\end{align*}
\begin{itemize}
\item The mean describes location; variance and standard deviation describe dispersion.
\item Standard deviation has the units of $Y$; variance has squared units.
\item For constants $a,b$: $\E[a+bY]=a+b\E[Y]$ and $\Var(a+bY)=b^2\Var(Y)$.
\end{itemize}
\begin{block}{Worked example}
If $Y$ takes $1,2,4$ with probabilities $.25,.50,.25$, then
$\E[Y]=2.25$, $\E[Y^2]=6.25$, and $\Var(Y)=6.25-2.25^2=1.1875$.
\end{block}
\end{frame}

\begin{frame}{Joint and Marginal Distributions: Read the Whole Table}
Let $X=1$ mean a sunny day and $Y=1$ mean buying ice cream.
\begin{center}
\begin{tabular}{c|cc|c}
 & $X=0$ & $X=1$ & Marginal of $Y$\\ \hline
$Y=0$ & .15 & .25 & .40\\
$Y=1$ & .05 & .55 & .60\\ \hline
Marginal of $X$ & .20 & .80 & 1
\end{tabular}
\end{center}
\begin{itemize}
\item A \textbf{joint probability} concerns both events: $\Prb(X=1,Y=1)=.55$.
\item A \textbf{marginal probability} ignores the other variable by summing over it:
\[
p_Y(y)=\sum_x p_{XY}(x,y),\qquad
p_X(x)=\sum_y p_{XY}(x,y).
\]
\item Thus $\Prb(Y=1)=.05+.55=.60$ and $\Prb(X=1)=.25+.55=.80$.
\end{itemize}
All joint cells are nonnegative and sum to 1. For continuous variables, integrate the joint density instead of summing probabilities.
\end{frame}

\begin{frame}{Conditional Distributions: Change the Reference Group}
For $\Prb(X=x)>0$,
\[
\Prb(Y=y\mid X=x)=\frac{\Prb(X=x,Y=y)}{\Prb(X=x)}.
\]
\begin{enumerate}
\item Restrict attention to outcomes with $X=x$.
\item Divide each relevant joint cell by the probability of that group.
\item Check that the conditional probabilities sum to 1.
\end{enumerate}
\begin{align*}
\Prb(Y=1\mid X=1)&=.55/.80=.6875,\\
\Prb(Y=0\mid X=1)&=.25/.80=.3125.
\end{align*}
Reversing the condition changes the denominator:
\[
\Prb(X=1\mid Y=1)=.55/.60\approx .9167.
\]
\textbf{Bayes' rule:}
$\Prb(X=x\mid Y=y)=\Prb(Y=y\mid X=x)\Prb(X=x)/\Prb(Y=y)$.
\end{frame}

\begin{frame}{Conditional Means and the Law of Iterated Expectations}
\small
The conditional mean averages $Y$ \textbf{within} a group:
\[
m(x)=\E[Y\mid X=x]=\sum_y y\,\Prb(Y=y\mid X=x).
\]
For binary $Y$, $\E[Y\mid X=x]=\Prb(Y=1\mid X=x)$.
In the sunny-day example,
\[
m(0)=.05/.20=.25,\qquad m(1)=.55/.80=.6875.
\]
The \textbf{Law of Iterated Expectations (LIE)} averages the group means using their population shares:
\begin{align*}
\E[Y]&=\E[\E[Y\mid X]]\\
&=\sum_x \E[Y\mid X=x]\Prb(X=x)\\
&=.25(.20)+.6875(.80)=.60.
\end{align*}
\textbf{What is random?} $m(x)$ is a number at a specified $x$; $m(X)=\E[Y\mid X]$ is a function of the random variable $X$.
Do not give all groups equal weight unless they have equal probability.
\end{frame}

\begin{frame}{Conditional-Expectation Rules You Will Use in Regression}
\small
\begin{enumerate}
\item \textbf{Linearity:}
$\E[aY+bZ\mid X]=a\E[Y\mid X]+b\E[Z\mid X]$ for constants $a,b$.
\item \textbf{Take out what is known:}
\[
\E[g(X)Y\mid X]=g(X)\E[Y\mid X].
\]
Why? Once $X=x$ is fixed, $g(x)$ is a constant:
\[
\E[g(X)Y\mid X=x]=g(x)\E[Y\mid X=x].
\]
In particular, $\E[XY\mid X]=X\E[Y\mid X]$.
\item \textbf{Conditional variance:}
\[
\Var(Y\mid X=x)=\E[(Y-m(x))^2\mid X=x].
\]
\end{enumerate}
A useful implication combines the first two rules with iterated expectations:
\[
\E[u\mid X]=0
\ \Longrightarrow\ 
\E[u]=0,\quad \E[Xu]=\E[X\E[u\mid X]]=0.
\]
Hence $\Cov(X,u)=0$.
\end{frame}

\begin{frame}{Independence, Covariance, and Correlation}
\small
\textbf{Independence:} knowing $X$ does not change the distribution of $Y$:
\[
p_{XY}(x,y)=p_X(x)p_Y(y)
\quad\Longleftrightarrow\quad
p_{Y\mid X}(y\mid x)=p_Y(y).
\]
\textbf{Covariance} measures linear comovement; \textbf{correlation} removes the units:
\[
\Cov(X,Y)=\E[XY]-\E[X]\E[Y],
\qquad
\rho_{XY}=\frac{\Cov(X,Y)}{\sigma_X\sigma_Y}\in[-1,1].
\]
\begin{itemize}
\item Independence implies zero covariance (with finite moments).
\item Zero covariance does \textbf{not} imply independence. If $X$ is equally likely to be $-1,0,1$ and $Y=X^2$, then $\Cov(X,Y)=0$ but $Y$ depends on $X$.
\item In the sunny-day table: $\Cov(X,Y)=.55-(.80)(.60)=.07>0$. The variables are not independent.
\end{itemize}
\[
\Var(aX+bY)=a^2\Var(X)+b^2\Var(Y)+2ab\Cov(X,Y).
\]
\end{frame}

\begin{frame}{Random Sampling and the Distribution of the Sample Mean}
An \textbf{independent and identically distributed (i.i.d.) sample} has:
\begin{itemize}
\item The same population distribution for every $Y_i$.
\item Independence across observations. This does not require $X_i$ and $Y_i$ within the same observation to be independent.
\end{itemize}
For i.i.d.\ observations with finite variance,
\begin{align*}
\overline Y&=\frac1n\sum_{i=1}^nY_i,\\
\E[\overline Y]&=\frac1n\sum_i\E[Y_i]=\mu_Y,\\
\Var(\overline Y)&=\frac1{n^2}\sum_i\Var(Y_i)
=\frac{n\sigma_Y^2}{n^2}=\frac{\sigma_Y^2}{n}.
\end{align*}
The covariance terms disappear because the observations are independent.
Thus $SD(\overline Y)=\sigma_Y/\sqrt n$: quadrupling $n$ halves this standard deviation.
\end{frame}

\begin{frame}{Law of Large Numbers: Why the Mean Is Consistent}
The \textbf{Law of Large Numbers (LLN)} describes convergence of the sample mean.
For i.i.d.\ observations with $\E[|Y_i|]<\infty$,
\[
\overline Y\xrightarrow{p}\mu_Y.
\]
Convergence in probability means that, for every tolerance $\varepsilon>0$,
\[
\Prb(|\overline Y-\mu_Y|>\varepsilon)\longrightarrow0.
\]
\begin{itemize}
\item As the sample grows, a large estimation error becomes unlikely.
\item This does \textbf{not} mean $\Prb(\overline Y=\mu_Y)\to1$. Exact equality need never occur for continuous data.
\item It does not mean that every additional observation moves the estimate closer to $\mu_Y$.
\end{itemize}
\begin{block}{Example}
If $Y_i$ indicates a head on an independent fair coin toss, $\overline Y$ is the sample fraction of heads and converges in probability to $0.5$.
\end{block}
\end{frame}

\begin{frame}{Central Limit Theorem: Why Normal Approximations Work}
The \textbf{Central Limit Theorem (CLT)} describes the shape of sampling uncertainty.
For i.i.d.\ observations with $0<\sigma_Y^2<\infty$,
\[
\frac{\overline Y-\mu_Y}{\sigma_Y/\sqrt n}
=\frac{\sqrt n(\overline Y-\mu_Y)}{\sigma_Y}
\xrightarrow{d}N(0,1).
\]
\begin{enumerate}
\item \textbf{Center:} subtract $\mu_Y$ to measure the estimation error.
\item \textbf{Scale:} divide by the standard deviation of the mean, $\sigma_Y/\sqrt n$.
\item \textbf{Approximate:} for sufficiently large $n$, compare this standardized error with a standard normal distribution.
\end{enumerate}
Equivalently, $\overline Y\approx N(\mu_Y,\sigma_Y^2/n)$ in large samples.
\begin{block}{LLN versus CLT}
The LLN says the error becomes small. The CLT describes its approximate distribution and scale. It does not say that individual observations become normal.
\end{block}
\end{frame}

\begin{frame}{Estimator Properties: Unbiasedness, Consistency, and Efficiency}
\small
\begin{center}
\begin{tabular}{p{2.7cm}p{4.4cm}p{6cm}}
\toprule
Property & Definition & Question it answers\\
\midrule
Unbiasedness & $\E[\widehat\theta]=\theta$ & Is the estimator centered on the truth across samples?\\[6pt]
Consistency & $\widehat\theta\xrightarrow{p}\theta$ & Does the estimation error vanish as $n$ grows?\\[6pt]
Efficiency & Smallest variance in a specified class & Which eligible estimator is most precise?\\
\bottomrule
\end{tabular}
\end{center}
\begin{itemize}
\item For i.i.d.\ data with finite variance, $\overline Y$ is unbiased and consistent.
\item Unbiasedness alone does not imply consistency: using $Y_1$ to estimate $\mu_Y$ is unbiased but its variance does not shrink with $n$.
\item When comparing biased estimators, variance alone is insufficient:
\[
MSE(\widehat\theta)=\Var(\widehat\theta)+\operatorname{Bias}(\widehat\theta)^2.
\]
Here $\operatorname{Bias}(\widehat\theta)=\E[\widehat\theta]-\theta$.
\end{itemize}
\end{frame}

\begin{frame}{BLUE, Least Squares, and the Cramér--Rao Bound}
\small
\textbf{BLUE} means \textbf{Best Linear Unbiased Estimator}.
For i.i.d.\ observations with common finite variance, a linear estimator $\sum_iw_iY_i$ is unbiased for every $\mu_Y$ if $\sum_iw_i=1$.
\[
\Var\!\left(\sum_iw_iY_i\right)
=\sigma_Y^2\sum_iw_i^2\geq\frac{\sigma_Y^2}{n}.
\]
Equal weights $w_i=1/n$ attain the bound: $\overline Y$ is BLUE.
\par\medskip
\textbf{Least squares} gives the same estimator by a different argument:
\[
\min_m\sum_i(Y_i-m)^2
\quad\Longrightarrow\quad
-2\sum_i(Y_i-m)=0
\quad\Longrightarrow\quad m=\overline Y.
\]
\textbf{Cramér--Rao:} under regularity conditions, an unbiased estimator satisfies
$\Var(\widehat\mu_Y)\geq1/\mathcal I_n(\mu_Y)$.
For normal observations with known variance, the sample's \textbf{Fisher information} is
$\mathcal I_n(\mu_Y)=n/\sigma_Y^2$. The mean reaches this bound even among nonlinear unbiased estimators.
\par\medskip
BLUE concerns a finite-sample comparison class; consistency is a separate large-sample property. Normality is not needed for the BLUE result above.
\end{frame}

\section{Lecture 2}
\begin{frame}{Estimating an Unknown Variance and the Mean's Standard Error}
\small
Before observing the sample, write the random sample variance as
\[
S_Y^2=\frac1{n-1}\sum_i(Y_i-\overline Y)^2.
\]
After observing it, compute the number
\[
s_Y^2=\frac1{n-1}\sum_i(y_i-\overline y)^2,\qquad s_Y=\sqrt{s_Y^2}.
\]
\begin{itemize}
\item We estimate $\mu_Y$ with $\overline y$. Since $\sum_i(y_i-\overline y)=0$, only $n-1$ deviations can vary freely.
\item Division by $n-1$ makes $S_Y^2$ unbiased for $\sigma_Y^2$ under i.i.d.\ sampling. Its square root need not be unbiased for $\sigma_Y$.
\item When $\sigma_Y$ is unknown, estimate the mean's SD by
\[
\boxed{SE(\overline Y)=\frac{s_Y}{\sqrt n}}.
\]
\end{itemize}
This guide writes $S_Y$ for the random statistic and $s_Y$ for its observed value. The lectures often use $s_Y$ for both; the distinction is the same as $\overline Y$ versus $\overline y$.
\end{frame}

\begin{frame}{Known $\sigma_Y$ or Estimated $s_Y$: Which Reference Distribution?}
\small
\begin{center}
\begin{tabular}{p{3.2cm}p{4.1cm}p{5.8cm}}
\toprule
Situation & Random test statistic under $H_0:\mu_Y=\mu_0$ & Reference distribution\\
\midrule
Known $\sigma_Y$ & $Z=\dfrac{\overline Y-\mu_0}{\sigma_Y/\sqrt n}$ & Exactly $N(0,1)$ for i.i.d.\ normal data; approximately normal for large $n$ by the CLT.\\[14pt]
Unknown $\sigma_Y$ & $T=\dfrac{\overline Y-\mu_0}{S_Y/\sqrt n}$ & Exactly Student's $t_{n-1}$ for i.i.d.\ normal data; approximately $N(0,1)$ for large $n$ under the CLT conditions.\\
\bottomrule
\end{tabular}
\end{center}
\begin{itemize}
\item The observed number is $t_{\mathrm{obs}}=(\overline y-\mu_0)/(s_Y/\sqrt n)$.
\item Estimating $\sigma_Y$ adds uncertainty: the finite-sample $t$ distribution has heavier tails.
\item In large samples, $S_Y/\sigma_Y\xrightarrow{p}1$, so replacing $\sigma_Y$ does not change the limiting standard normal distribution.
\end{itemize}
For small, non-normal samples, neither an exact $t$ claim nor an automatic normal approximation is justified.
\end{frame}

\begin{frame}{Hypothesis Testing: A Procedure, Not Just a Formula}
\begin{enumerate}
\item \textbf{State the parameter and null:} for example, $H_0:\mu_Y=175$ cm.
\item \textbf{Choose the alternative:} $H_1:\mu_Y\neq175$ is two-sided; $H_1:\mu_Y>175$ is one-sided.
\item \textbf{Choose a significance level} $\alpha$ before looking at the result.
\item \textbf{Compute the standard error} and the observed statistic:
\[
t_{\mathrm{obs}}=\frac{\overline y-\mu_0}{SE(\overline Y)}.
\]
\item \textbf{Use the appropriate null distribution} to calculate a $p$-value.
\item \textbf{Decide and interpret:} reject if $p\leq\alpha$; otherwise do not reject.
\end{enumerate}
\begin{block}{Interpret the statistic}
A value $t_{\mathrm{obs}}=-2.4$ means the observed mean is 2.4 estimated standard errors below the null value. It is not a probability.
\end{block}
\end{frame}

\begin{frame}{$p$-Values, Tail Areas, and Significance Levels}
\small
A \textbf{$p$-value is a probability under the null hypothesis}: how often would the test statistic be at least as extreme as the observed one?
For $Z\sim N(0,1)$, define $\Phi(z)=\Prb(Z\leq z)$.
\begin{center}
\begin{tabular}{lll}
\toprule
Alternative & Normal-approximation $p$-value & Reject at 5\% if\\
\midrule
$\mu_Y\neq\mu_0$ & $2[1-\Phi(|t_{\mathrm{obs}}|)]$ & $|t_{\mathrm{obs}}|\geq1.96$\\
\bottomrule
\end{tabular}
\end{center}
For two-sided normal tests, critical values are:
\[
\alpha=10\%:\ 1.645,\qquad
\alpha=5\%:\ 1.960,\qquad
\alpha=1\%:\ 2.576.
\]
Use Student's $t$ probabilities and critical values when that is the reference distribution.
\par\medskip
\textbf{Avoid these interpretations:} the $p$-value is not $\Prb(H_0\text{ is true})$; failure to reject does not establish that $H_0$ is true. Statistical significance does not measure economic importance.
\end{frame}

\begin{frame}{Worked Example: Height, Standard Error, and a Mean Test}
\small
The ten illustrative heights from Lecture 2 have
\[
n=10,\quad \overline y=170\text{ cm},\quad
\sum_i(y_i-\overline y)^2=388\text{ cm}^2.
\]
\begin{align*}
s_Y&=\sqrt{388/9}=6.5659\text{ cm},\\
SE(\overline Y)&=6.5659/\sqrt{10}=2.0763\text{ cm},\\
t_{\mathrm{obs}}&=\frac{170-175}{2.0763}=-2.4081.
\end{align*}
For $H_0:\mu_Y=175$ versus $H_1:\mu_Y\neq175$:
\begin{itemize}
\item \textbf{Normal approximation:} $p\approx .0160$ and $|t_{\mathrm{obs}}|>1.96$.
\item \textbf{If heights are i.i.d.\ normal and $\sigma_Y$ is unknown:} use $t_9$; $p\approx .0394$ and the 5\% critical value is $2.262$.
\item Both reject at 5\%, but their $p$-values differ. With only ten observations, normality is an additional assumption, not a consequence of a large sample.
\end{itemize}
Do not divide by $s_Y$ alone: that describes individual dispersion, not uncertainty in the mean.
\end{frame}

\begin{frame}{Confidence Intervals: Formula, Meaning, and Link to Tests}
\small
A two-sided $(1-\alpha)$ confidence interval for a mean is
\[
\overline y\ \pm\ c\,SE(\overline Y),
\]
where $c=z_{1-\alpha/2}$ for a normal approximation, or $c=t_{1-\alpha/2,n-1}$ for exact inference with i.i.d.\ normal data and unknown variance.
\begin{itemize}
\item \textbf{Coverage:} over repeated samples, the interval procedure contains the fixed true parameter in $100(1-\alpha)\%$ of samples (approximately for asymptotic intervals).
\item \textbf{Width:} more confidence gives a wider interval; greater precision gives a narrower interval.
\item \textbf{Test equivalence:} a null value outside the interval is rejected by the corresponding two-sided test, using the same standard error and reference distribution.
\end{itemize}
For the height example:
\[
\text{Normal 95\% CI: }170\pm1.96(2.0763)\approx[165.93,174.07],
\]
\[
\text{Normal-population }t_9\text{ CI: }170\pm2.262(2.0763)\approx[165.30,174.70].
\]
These are intervals for the \textbf{population mean}, not ranges containing 95\% of individual heights.
\end{frame}

\begin{frame}{Type I Error, Type II Error, and Power}
\begin{center}
\begin{tabular}{lcc}
\toprule
 & $H_0$ true & A particular alternative true\\
\midrule
Reject $H_0$ & Type I error & Correct rejection\\
Do not reject $H_0$ & Correct non-rejection & Type II error\\
\bottomrule
\end{tabular}
\end{center}
For a test calibrated to have size $\alpha$,
\[
\alpha=\Prb(\text{reject }H_0\mid H_0\text{ true}).
\]
At a specified alternative mean $\mu_1$,
\[
\beta(\mu_1)=\Prb_{\mu_1}(\text{do not reject }H_0),
\qquad
\operatorname{Power}(\mu_1)=1-\beta(\mu_1).
\]
\begin{itemize}
\item $\alpha$ controls false rejection under the null. $\beta$ depends on the particular true alternative, sample size, and noise.
\item For a fixed sample size and test family, lowering $\alpha$ makes rejection harder and generally reduces power.
\item Larger samples or effects farther from the null generally increase power.
\end{itemize}
\end{frame}

\begin{frame}{How to Calculate $\beta$: Use a Specific Alternative}
\small
Consider a two-sided normal mean test with known $\sigma_Y$ and rejection rule
\[
|\overline Y-\mu_0|>c\,\frac{\sigma_Y}{\sqrt n},
\qquad c=z_{1-\alpha/2}.
\]
If the true mean is $\mu_1$, define the effect in standard-error units:
\[
\delta=\frac{\mu_1-\mu_0}{\sigma_Y/\sqrt n}.
\]
Under this alternative, the test statistic has distribution $N(\delta,1)$, so
\[
\boxed{\beta(\mu_1)=\Phi(c-\delta)-\Phi(-c-\delta)}.
\]
\begin{block}{Example}
Let $\alpha=.05$, $c=1.96$, and $\delta=2$.
Then $\beta=\Phi(-.04)-\Phi(-3.96)\approx .4840$, so power is about $.5160$.
\end{block}
The formula is exact for normal observations with known variance and approximate when justified by the CLT. There is no single $\beta$ determined by $\alpha$ alone.
\end{frame}

\begin{frame}{Comparing Two Means: Independent Groups}
\small
For independent samples from groups $a$ and $b$, let $\Delta=\mu_a-\mu_b$ and test $H_0:\Delta=d_0$.
\begin{align*}
\widehat\Delta&=\overline Y^a-\overline Y^b,\\
\Var(\widehat\Delta)&=\frac{\sigma_a^2}{n_a}+\frac{\sigma_b^2}{n_b},\\
SE(\widehat\Delta)&=\sqrt{\frac{s_a^2}{n_a}+\frac{s_b^2}{n_b}},\\
t_{\mathrm{obs}}&=\frac{(\overline y^a-\overline y^b)-d_0}{SE(\widehat\Delta)}.
\end{align*}
\begin{itemize}
\item Subtract the hypothesized difference $d_0$, which need not be zero.
\item Large-sample 95\% CI: $(\overline y^a-\overline y^b)\pm1.96\,SE(\widehat\Delta)$.
\item Independence removes the covariance term. For paired measurements on the same people, analyze within-person differences instead.
\item For small independent normal samples with unknown unequal variances, use Welch's $t$ method rather than treating 1.96 as exact.
\end{itemize}
\end{frame}

\begin{frame}{Group Differences: Calculation and Causal Interpretation}
\small
From the two height samples in Lecture 2:
\[
n_a=n_b=10,\quad \overline y^a=170,\quad\overline y^b=164,
\quad s_a^2=388/9,\quad s_b^2=132/9.
\]
\[
SE(\widehat\Delta)=\sqrt{\frac{388/9}{10}+\frac{132/9}{10}}=2.4037.
\]
For $H_0:\Delta=5$ cm:
\[
t_{\mathrm{obs}}=\frac{(170-164)-5}{2.4037}=.4160,
\qquad p\approx .6774
\]
using the lecture's rough normal approximation. Do not reject at 5\%.
\par\medskip
\textbf{A difference in conditional means is initially an association:}
\[
\E[Y\mid D=1]-\E[Y\mid D=0].
\]
Random assignment, together with a well-defined treatment, no spillovers, and appropriate observation of outcomes, can justify a causal interpretation.
In observational data, different group composition can explain part or all of the difference.
\end{frame}

\section{Lecture 3}
\begin{frame}{Simple Linear Regression: Three Different Objects}
\small
\begin{center}
\begin{tabular}{>{\raggedright\arraybackslash}p{3.2cm}>{\raggedright\arraybackslash}p{5.4cm}>{\raggedright\arraybackslash}p{4.6cm}}
\toprule
Object & Equation & Meaning\\
\midrule
Population model & $Y_i=\beta_0+\beta_1X_i+u_i$ & Unknown parameters and unobserved error\\[9pt]
Estimated rule before observing data & $\widehat Y_i=\widehat\beta_0+\widehat\beta_1X_i$ & Random estimators and fitted values\\[9pt]
Fitted line in an observed sample & $\widehat y_i=\widehat b_0+\widehat b_1x_i$ & Numerical estimates and predictions\\
\bottomrule
\end{tabular}
\end{center}
\begin{itemize}
\item $\beta_0$ is an intercept; $\beta_1$ is a change in $Y$ per unit of $X$.
\item $u_i$ contains determinants of $Y_i$ omitted from the displayed part of the model.
\item A sample residual $\widehat u_i$ is not the same object as the population error $u_i$.
\end{itemize}
For clarity, this guide uses $\widehat b_j$ for a realized estimate of $\widehat\beta_j$; the lectures often retain $\widehat\beta_j$ in numerical applications.
\end{frame}

\begin{frame}{Ordinary Least Squares: Objective and First-Order Conditions}
\small
\textbf{Ordinary least squares (OLS)} chooses an intercept $a$ and slope $b$ to minimize
\[
Q_n(a,b)=\sum_{i=1}^n(Y_i-a-bX_i)^2.
\]
The \textbf{first-order conditions (FOCs)} are
\begin{align*}
\frac{\partial Q_n}{\partial a}
&=-2\sum_i(Y_i-a-bX_i)=0,\\
\frac{\partial Q_n}{\partial b}
&=-2\sum_iX_i(Y_i-a-bX_i)=0.
\end{align*}
Evaluate them at $(a,b)=(\widehat\beta_0,\widehat\beta_1)$.
The intercept condition gives
\[
n\overline Y-n\widehat\beta_0-n\widehat\beta_1\overline X=0
\quad\Longrightarrow\quad
\boxed{\widehat\beta_0=\overline Y-\widehat\beta_1\overline X}.
\]
If $TSS_X=\sum_i(X_i-\overline X)^2>0$, the squared-error objective has a unique minimizer.
\end{frame}

\begin{frame}{Solve the Slope FOC: Center, Substitute, and Divide}
\small
Substitute the intercept expression into the slope FOC:
\begin{align*}
0&=\sum_iX_i[Y_i-(\overline Y-\widehat\beta_1\overline X)-\widehat\beta_1X_i]\\
 &=\sum_iX_i[(Y_i-\overline Y)-\widehat\beta_1(X_i-\overline X)].
\end{align*}
Since $\sum_i(Y_i-\overline Y)=0$ and $\sum_i(X_i-\overline X)=0$,
\[
\sum_iX_i(Y_i-\overline Y)
=\sum_i(X_i-\overline X)(Y_i-\overline Y),
\]
\[
\sum_iX_i(X_i-\overline X)=\sum_i(X_i-\overline X)^2=TSS_X.
\]
$TSS_X$ is the \textbf{Total Sum of Squares of $X$}: the total variation of the sampled $X_i$ values around their sample mean.
Therefore,
\[
\boxed{\widehat\beta_1=
\frac{\sum_i(X_i-\overline X)(Y_i-\overline Y)}{TSS_X}
=\frac{S_{XY}}{S_X^2}}.
\]
Here $S_{XY}$ and $S_X^2$ denote the random sample covariance and variance, each dividing its sum by $n-1$. Their observed values are $s_{XY}$ and $s_X^2$.
\end{frame}

\begin{frame}{Worked OLS Example: From Sums to a Fitted Line}
\small
For the ten observations in Lecture 3, the observed summaries are
\[
n=10,\quad \sum_i x_i=55,\quad \sum_i y_i=84.90,
\quad \overline x=5.5,\quad \overline y=8.49,
\]
\[
\sum_i(x_i-\overline x)^2=82.50,\qquad
\sum_i(x_i-\overline x)(y_i-\overline y)=-31.75.
\]
\begin{enumerate}
\item Compute the sample variance and covariance:
$s_X^2=82.50/9=9.1667$, $s_{XY}=-31.75/9=-3.5278$.
\item Divide to obtain the slope:
$\widehat b_1=-31.75/82.50=-.38485$.
\item Recover the intercept:
$\widehat b_0=8.49-(-.38485)(5.5)=10.60667$.
\end{enumerate}
Thus $\widehat y_i=10.60667-.38485x_i$.
At $x_1=1$, $\widehat y_1=10.22182$; since $y_1=10.8$, the residual is $.57818$.
\par\medskip
A one-unit increase in $X$ is associated with about $.385$ fewer units of $Y$. This slope is not automatically causal.
\end{frame}

\begin{frame}{OLS Identities: What Is Zero by Construction?}
\small
For $Q_n(a,b)=\sum_i(Y_i-a-bX_i)^2$, evaluate the first-order conditions at $(a,b)=(\widehat\beta_0,\widehat\beta_1)$:
\begin{align*}
\left.\frac{\partial Q_n}{\partial a}\right|_{\widehat\beta_0,\widehat\beta_1} &=-2\sum_i\widehat u_i=0 \quad\Longrightarrow\quad \sum_i\widehat u_i=0,\\
\left.\frac{\partial Q_n}{\partial b}\right|_{\widehat\beta_0,\widehat\beta_1} &=-2\sum_iX_i\widehat u_i=0 \quad\Longrightarrow\quad \sum_iX_i\widehat u_i=0.
\end{align*}
The first is the intercept FOC; the second is the slope FOC.
\begin{itemize}
\item \textbf{Mean residual is zero:} $\overline{\widehat u}=0$. Individual residuals are usually not zero.
\item \textbf{The line passes through the means:}
$\overline Y=\widehat\beta_0+\widehat\beta_1\overline X$.
\item \textbf{Fitted values and residuals are orthogonal:}
\[
\sum_i\widehat Y_i\widehat u_i
=\widehat\beta_0\sum_i\widehat u_i+
\widehat\beta_1\sum_iX_i\widehat u_i=0.
\]
\end{itemize}
\begin{block}{An algebraic identity is not an identifying assumption}
These identities hold for the fitted OLS residuals even when $\E[u_i\mid X_i]\neq0$.
A zero sample residual mean does not prove conditional mean zero for the population error.
\end{block}
\end{frame}

\begin{frame}{Goodness of Fit: Three Sums of Squares and $R^2$}
\small
\begin{align*}
TSS_Y&=\sum_i(Y_i-\overline Y)^2 &&\text{total variation around the mean},\\
ESS&=\sum_i(\widehat Y_i-\overline Y)^2 &&\text{variation explained by the fitted line},\\
SSR&=\sum_i(Y_i-\widehat Y_i)^2 &&\text{residual sum of squares}.
\end{align*}
Because $Y_i-\overline Y=(\widehat Y_i-\overline Y)+\widehat u_i$ and the cross term sums to zero,
\[
TSS_Y=ESS+SSR,\qquad
\boxed{R^2=\frac{ESS}{TSS_Y}=1-\frac{SSR}{TSS_Y}}.
\]
\begin{itemize}
\item With an intercept and $TSS_Y>0$, $0\leq R^2\leq1$.
\item $R^2=.60$ means the fitted line accounts for 60\% of the sample variation around the mean.
\item High $R^2$ does not establish causality or guarantee good predictions in a new sample.
\end{itemize}
Keep $TSS_X$ (variation in $X$) separate from $TSS_Y$ (variation in $Y$).
\end{frame}

\begin{frame}{SER, RMSE, and Coefficient Standard Errors Are Different}
\small
\begin{columns}[T]
\begin{column}{.48\textwidth}
\begin{block}{Standard error of the regression (SER)}
\[
\widehat\sigma_u^2=\frac{SSR}{n-2}=SER^2,
\qquad SER=\sqrt{\widehat\sigma_u^2}.
\]
Two estimated coefficients use two degrees of freedom. Under the homoskedastic model, $SER^2$ estimates $\sigma_u^2$.
\end{block}
\end{column}
\begin{column}{.48\textwidth}
\begin{block}{Root mean squared error (RMSE)}
\[
RMSE=\sqrt{\frac{SSR}{n}}.
\]
The square root of the average squared in-sample prediction error:
$Y_i-\widehat Y_i=\widehat u_i$.
\end{block}
\end{column}
\end{columns}
\medskip
\begin{itemize}
\item Both measure unexplained outcome dispersion, in units of $Y$; they become similar when $n$ is large.
\item $SE(\widehat\beta_1)$ instead measures uncertainty about the estimated slope, in units of $Y$ per unit of $X$.
\item Neither $SER$ nor $\widehat Y\pm SER$ is by itself a calibrated 95\% prediction interval.
\end{itemize}
Always check the denominator: some software labels $\sqrt{SSR/(n-2)}$ as Root MSE.
\end{frame}

\begin{frame}{Estimating the Error Variance: Why Residuals Need a Correction}
\small
\textbf{Assumptions.} Condition on all sampled regressors $\mathbf X$:
$\E[u_i\mid\mathbf X]=0$, $\Var(u_i\mid\mathbf X)=\sigma_u^2$, and errors are
uncorrelated conditional on $\mathbf X$; $TSS_X>0$ and $n>2$.
\begin{block}{If the population coefficients were known}
We could observe $u_i=Y_i-\beta_0-\beta_1X_i$. Since
$\E[u_i^2\mid\mathbf X]=\sigma_u^2$,
\[
\E\!\left[\frac1n\sum_i u_i^2\,\middle|\,\mathbf X\right]=\sigma_u^2.
\]
\end{block}
Actual residuals use fitted coefficients:
$\widehat u_i=Y_i-\widehat\beta_0-\widehat\beta_1X_i$.
Least squares fits some of the realized error, so the residual sum of squares is smaller on average.
\begin{align*}
\E[SSR\mid\mathbf X]&=(n-2)\sigma_u^2
&&\text{(derived on the next slide)},\\
\widehat\sigma_u^2&=\frac{SSR}{n-2},
&\E[\widehat\sigma_u^2\mid\mathbf X]&=\sigma_u^2.
\end{align*}
\textbf{SER estimates $\sigma_u$; $SER^2$ is the unbiased variance estimator under these assumptions.}
\end{frame}

\begin{frame}{Why the Correction Is $n-2$: A Derivation}
\small
Under the previous assumptions, define $\overline u=n^{-1}\sum_i u_i$ and
$d=\widehat\beta_1-\beta_1$. The intercept identity gives
\[
\widehat\beta_0-\beta_0=\overline u-d\overline X
\quad\Longrightarrow\quad
u_i=\widehat u_i+\overline u+d(X_i-\overline X).
\]
OLS orthogonality, $\sum_i\widehat u_i=0$ and
$\sum_i(X_i-\overline X)\widehat u_i=0$, together with
$\sum_i(X_i-\overline X)=0$, eliminates all cross-products when squaring and summing:
\[
SSR=\sum_i u_i^2-n\overline u^{\,2}-d^2TSS_X.
\]
Using OLS unbiasedness and $\Var(\widehat\beta_1\mid\mathbf X)=\sigma_u^2/TSS_X$
\textbf{(results derived later in this guide)}, we obtain
\begin{align*}
\E\!\left[\sum_i u_i^2\,\middle|\,\mathbf X\right]&=n\sigma_u^2,\\
\E[n\overline u^{\,2}\mid\mathbf X]&=n\Var(\overline u\mid\mathbf X)=\sigma_u^2,\\
\E[d^2TSS_X\mid\mathbf X]&=TSS_X\Var(\widehat\beta_1\mid\mathbf X)=\sigma_u^2.
\end{align*}
Thus $\E[SSR\mid\mathbf X]=n\sigma_u^2-\sigma_u^2-\sigma_u^2=(n-2)\sigma_u^2$.
One fitted direction for the intercept and one for the slope remove two $\sigma_u^2$ terms.
\end{frame}

\begin{frame}{Population Least Squares and Sample OLS}
\small
The \textbf{population linear projection} minimizes an expected squared error:
\[
(a^*,b^*)=\arg\min_{a,b}\E[(Y-a-bX)^2].
\]
Its FOCs, when the moments exist, are
\[
\E[Y-a^*-b^*X]=0,\qquad
\E[X(Y-a^*-b^*X)]=0.
\]
Thus, for $\Var(X)>0$,
\[
b^*=\frac{\Cov(X,Y)}{\Var(X)},\qquad
a^*=\E[Y]-b^*\E[X].
\]
\begin{itemize}
\item Population coefficients are fixed features of the joint distribution.
\item Sample OLS replaces population moments with sample moments; the estimated coefficients vary across samples.
\item If $\E[Y\mid X]=\beta_0+\beta_1X$, the projection equals that conditional-mean line.
\item If the conditional mean is nonlinear, the projection remains the best \emph{linear} predictor under squared loss.
\end{itemize}
Here $a^*,b^*$ distinguish population projection coefficients from sample estimators and any structural causal parameters.
\end{frame}

\begin{frame}{OLS Assumptions: What Each One Does}
\small
For a model $Y_i=\beta_0+\beta_1X_i+u_i$:
\begin{enumerate}
\item \textbf{Conditional mean zero:} $\E[u_i\mid X_i]=0$.
It implies $\E[Y_i\mid X_i]=\beta_0+\beta_1X_i$ and removes systematic error differences across values of $X_i$.
\item \textbf{i.i.d.\ observations:} the pairs $(X_i,Y_i)$ are independent across $i$ and drawn from one joint distribution.
This supports the sampling arguments.
\item \textbf{Finite moments and regressor variation:} the lecture assumes finite fourth moments and $\Var(X)>0$.
These control tail behavior and ensure that a slope can be identified and estimated.
\end{enumerate}
\begin{block}{Causal interpretation needs a causal model}
If $\beta_1$ is a structural causal effect, conditional mean zero helps identify it from observed variation. A fitted conditional-mean relationship alone does not prove that intervening on $X$ changes $Y$.
\end{block}
\textbf{Homoskedasticity}, $\Var(u_i\mid X_i)=\sigma_u^2$, is an additional assumption for the simple variance formula; it is not required for OLS unbiasedness.
\end{frame}

\begin{frame}{Omitted Variables: Why Conditional Mean Zero Can Fail}
\small
Suppose the structural equation is
\[
Y_i=\beta_0+\beta_1X_i+\gamma Z_i+v_i,\qquad \E[v_i\mid X_i,Z_i]=0.
\]
If we omit $Z_i$ and keep the displayed coefficients $\beta_0,\beta_1$, comparison with
$Y_i=\beta_0+\beta_1X_i+u_i$ gives
\[
u_i=\gamma Z_i+v_i,\qquad
\E[u_i\mid X_i]=\gamma\E[Z_i\mid X_i].
\]
If $\gamma\neq0$ and the conditional mean of $Z_i$ varies with $X_i$, the conditional mean of $u_i$ varies too; an intercept shift cannot remove that variation.
\par\medskip
The population projection slope becomes
\[
b^*=\frac{\Cov(X,Y)}{\Var(X)}
=\beta_1+\gamma\frac{\Cov(X,Z)}{\Var(X)}.
\]
For example, if $Z$ is disadvantage, $\gamma<0$, and $\Cov(X,Z)>0$, omitting it shifts the slope downward.
Conditional mean zero implies zero covariance; its failure does not always imply nonzero covariance.
\end{frame}

\begin{frame}{Decompose the OLS Slope: True Coefficient Plus Sampling Error}
\small
Write $D_i=X_i-\overline X$ and $TSS_X=\sum_iD_i^2$. First,
\[
\sum_iD_i(Y_i-\overline Y)=\sum_iD_iY_i
\quad\text{because}\quad \sum_iD_i=0.
\]
Substitute $Y_i=\beta_0+\beta_1X_i+u_i$:
\[
\widehat\beta_1
=\frac{\beta_0\sum_iD_i+\beta_1\sum_iD_iX_i+\sum_iD_i u_i}{TSS_X}.
\]
Now use $X_i=\overline X+D_i$:
\[
\sum_iD_iX_i=\overline X\sum_iD_i+\sum_iD_i^2=TSS_X.
\]
Therefore,
\[
\boxed{\widehat\beta_1=\beta_1+\frac{\sum_i(X_i-\overline X)u_i}{TSS_X}}.
\]
The intercept term disappears, and the remaining fraction is the estimation error. This identity holds before taking expectations.
\end{frame}

\begin{frame}{Why OLS Is Unbiased: Condition First, Then Average}
\small
Let $\mathbf X=(X_1,\ldots,X_n)$ denote \textbf{all} sampled regressors.
Under i.i.d.\ sampling and conditional mean zero, $\E[u_i\mid\mathbf X]=0$.
\begin{enumerate}
\item Conditional on $\mathbf X$, every $X_i$, $\overline X$, and $TSS_X$ is fixed.
\item Take the conditional expectation of the decomposition:
\[
\E[\widehat\beta_1\mid\mathbf X]
=\beta_1+\frac{\sum_i(X_i-\overline X)\E[u_i\mid\mathbf X]}{TSS_X}
=\beta_1.
\]
\item Apply the Law of Iterated Expectations:
\[
\E[\widehat\beta_1]=\E[\E[\widehat\beta_1\mid\mathbf X]]=\beta_1.
\]
\end{enumerate}
The next slide derives the corresponding result for the intercept.
\par\medskip
These results require an identified slope ($TSS_X>0$) and existence of the relevant expectations. Homoskedasticity and normality are not needed for this argument.
\end{frame}

\begin{frame}{Why the Intercept Is Unbiased: Step by Step}
\small
Define $\overline u=\frac1n\sum_i u_i$.
\begin{enumerate}
\item Average the population equation:
\[
\overline Y=\beta_0+\beta_1\overline X+\overline u.
\]
\item Substitute into the intercept estimator:
\[
\widehat\beta_0=\overline Y-\widehat\beta_1\overline X
=\beta_0+(\beta_1-\widehat\beta_1)\overline X+\overline u.
\]
\item Conditional on $\mathbf X=(X_1,\ldots,X_n)$, $\overline X$ is fixed and
\[
\E[\overline u\mid\mathbf X]=\frac1n\sum_i\E[u_i\mid\mathbf X]=0.
\]
Using the previous result $\E[\widehat\beta_1\mid\mathbf X]=\beta_1$,
\[
\begin{aligned}
\E[\widehat\beta_0\mid\mathbf X]
&=\beta_0+(\beta_1-\E[\widehat\beta_1\mid\mathbf X])\overline X
+\E[\overline u\mid\mathbf X]\\
&=\beta_0+(\beta_1-\beta_1)\overline X+0=\beta_0.
\end{aligned}
\]
\item Apply iterated expectations:
\[
\E[\widehat\beta_0]=\E[\E[\widehat\beta_0\mid\mathbf X]]=\beta_0.
\]
\end{enumerate}
\end{frame}

\begin{frame}{OLS Precision: Variance and Standard Error of the Slope}
\small
With errors uncorrelated across observations conditional on $\mathbf X$,
\begin{align*}
\Var(\widehat\beta_1\mid\mathbf X)
&=\frac{1}{TSS_X^2}\Var\!\left(\sum_i(X_i-\overline X)u_i\mid\mathbf X\right)\\
&=\frac{\sum_i(X_i-\overline X)^2\Var(u_i\mid\mathbf X)}{TSS_X^2}.
\end{align*}
Under \textbf{homoskedasticity}, every conditional error variance is $\sigma_u^2$, so
\[
\Var(\widehat\beta_1\mid\mathbf X)=\frac{\sigma_u^2}{TSS_X},
\qquad
\boxed{SE(\widehat\beta_1)=\sqrt{\frac{SSR}{(n-2)TSS_X}}
=\frac{SER}{\sqrt{TSS_X}}}.
\]
\begin{itemize}
\item More error variation makes the slope less precise.
\item More variation in $X$ increases $TSS_X$ and makes the slope more precise.
\item With stable regressor variation, larger samples typically increase $TSS_X$.
\end{itemize}
If error variance varies with $X$, the simple homoskedastic standard-error formula need not be valid.
\end{frame}

\begin{frame}{Precision of the Intercept: Where Is $X=0$?}
\small
Under conditional mean zero and homoskedastic errors that are uncorrelated across observations conditional on $\mathbf X$,
\[
\Var(\widehat\beta_0\mid\mathbf X)
=\sigma_u^2\left(\frac1n+\frac{\overline X^2}{TSS_X}\right).
\]
Replace the unknown error variance by $SER^2=SSR/(n-2)$ to estimate its standard error:
\[
\boxed{SE(\widehat\beta_0)
=SER\sqrt{\frac1n+\frac{\overline X^2}{TSS_X}}}.
\]
\begin{itemize}
\item The intercept is the fitted outcome at $X=0$.
\item Holding $n$, error variance, and $TSS_X$ fixed, the intercept is less precise when the sample's $X$ values are centered farther from zero.
\item If $\overline X=0$, then $SE(\widehat\beta_0)=SER/\sqrt n$.
\item A small standard error does not make an extrapolation substantively meaningful: check whether $X=0$ is relevant to the population.
\end{itemize}
Use the standard error for the coefficient being studied; the intercept and slope generally have different standard errors and different units.
\end{frame}

\begin{frame}{OLS in Large Samples: Consistency and Normal Approximation}
\small
Divide the numerator and denominator of the estimation error by $n$:
\[
\widehat\beta_1-\beta_1
=\frac{n^{-1}\sum_i(X_i-\overline X)u_i}
       {n^{-1}\sum_i(X_i-\overline X)^2}.
\]
\begin{enumerate}
\item By the \textbf{Law of Large Numbers}, the numerator converges to $\Cov(X,u)=0$ under conditional mean zero.
\item The denominator converges to $\Var(X)>0$.
\item Therefore $\widehat\beta_1\xrightarrow{p}\beta_1$: OLS is \textbf{consistent}.
\end{enumerate}
Under the lecture's i.i.d.\ and moment assumptions, the \textbf{Central Limit Theorem} also gives asymptotic normality.
With homoskedasticity,
\[
\sqrt n(\widehat\beta_1-\beta_1)
\xrightarrow{d}N\!\left(0,\frac{\sigma_u^2}{\Var(X)}\right).
\]
Dividing $\widehat\beta_1-\beta_1$ by an appropriate estimated standard error gives an approximately standard normal statistic in large samples. This parallels inference for $\overline Y$.
\end{frame}

\section{Review}
\begin{frame}{Lecture 1: Questions You Should Be Able to Answer}
\small
\begin{enumerate}
\item \textbf{Joint, marginal, or conditional?}
In the sunny-day table, explain why $.55$, $.60$, and $.6875$ answer different questions.
\item \textbf{Why not reverse the condition?}
Explain why $\Prb(Y=1\mid X=1)$ and $\Prb(X=1\mid Y=1)$ have different denominators.
\item \textbf{Which average?}
Recover $\E[Y]=.60$ by weighting the two conditional means.
\item \textbf{What varies across samples?}
Distinguish $\mu_Y$, $Y_i$, $\overline Y$, and $\overline y$.
\item \textbf{Which theorem?}
Use the Law of Large Numbers to explain consistency and the Central Limit Theorem to explain a normal approximation.
\item \textbf{Which implication is valid?}
Independence implies zero covariance; the converse is false.
\end{enumerate}
\begin{block}{Key check}
If $n$ increases from 25 to 100 with the same population variance, the SD of the sample mean halves. The population SD of individual observations stays the same.
\end{block}
\end{frame}

\begin{frame}{Lecture 2: An Inference Checklist}
\begin{enumerate}
\item Identify the parameter: one mean or a difference of means.
\item State $H_0$, $H_1$, and $\alpha$; determine whether the test is one- or two-sided.
\item Determine whether $\sigma_Y$ is known. If not, estimate it with $s_Y$.
\item Calculate the \textbf{standard error of the estimator}, not just the SD of the observations.
\item Compute the observed statistic and choose a justified reference distribution.
\item Calculate the $p$-value or compare with a critical value; interpret the decision.
\item Construct a confidence interval and check consistency with the corresponding test.
\end{enumerate}
\begin{block}{Quick check}
If $\overline y=52$, $\mu_0=50$, $s_Y=10$, and $n=100$, then
$SE(\overline Y)=1$, $t_{\mathrm{obs}}=2$, and the two-sided normal-approximation $p$-value is about $.0455$. The 95\% interval is $[50.04,53.96]$.
\end{block}
\end{frame}

\begin{frame}{Lecture 3: An OLS Checklist}
\small
\begin{enumerate}
\item Compute $\overline x,\overline y$, squared deviations of $x$, and cross-products.
\item Obtain $\widehat b_1=s_{XY}/s_X^2$ and $\widehat b_0=\overline y-\widehat b_1\overline x$.
\item Calculate $\widehat y_i$ and residuals $y_i-\widehat y_i$. Check the two sample FOCs.
\item Compute $SSR,TSS_Y,ESS,R^2$, and $SER$; interpret their units and meanings.
\item State the assumptions needed for unbiasedness, consistency, and the standard-error formula.
\item Explain whether the slope has a causal interpretation; identify plausible omitted variables.
\end{enumerate}
\begin{block}{Do not confuse these statements}
$\sum_i\widehat u_i=0$ is an OLS identity.\\
$\E[u_i\mid X_i]=0$ is a population assumption.\\
A high $R^2$ describes fit. A small standard error describes precision.\\
Neither one alone establishes causality.
\end{block}
\end{frame}

\end{document}
